\documentclass[11pt]{article}

\usepackage[margin=1in]{geometry}
\usepackage{amsmath,amssymb,amsfonts}
\usepackage{booktabs}
\usepackage{enumitem}
\usepackage{hyperref}
\usepackage{xcolor}

\title{DUC-WM: LLM-Guided Mechanism-Residual World Models for Planning under Hidden Dynamics}
\author{Draft for mentor discussion}
\date{\today}

\newcommand{\R}{\mathbb{R}}
\newcommand{\E}{\mathbb{E}}
\newcommand{\KL}{\mathrm{KL}}
\newcommand{\diag}{\mathrm{diag}}
\newcommand{\clip}{\mathrm{clip}}
\newcommand{\softplus}{\mathrm{softplus}}

\begin{document}
\maketitle

\begin{abstract}
This note describes DUC-WM, a model-based control framework for continuous-control environments with hidden mechanism shifts such as wind, friction, mass, damping, delay, sticky transitions, impulses, and gravity changes. The central idea is to use an LLM only as a constrained scientific prior builder: it proposes a small set of named mechanisms, sparse input/output masks, prior confidence, and a safe law type from a fixed DSL. These priors are compiled into executable tensor laws, validated against observed transition deltas, and then repaired by mechanism-specific neural residuals. A history encoder infers mechanism strengths from recent trajectories, so deployment does not require oracle context labels. The learned world model is evaluated not only by rollout prediction accuracy, but also by its utility for CEM-MPC planning.

The framework contribution is not that LLM priors, residual world models, or latent-context models are individually new. The novelty is their controlled integration into a mechanism-residual world model with executable priors, data-validated prior gates, residual trust regions, hidden-context inference, and a fair planning-oriented evaluation protocol.
\end{abstract}

\section{Problem Setting}

We consider continuous-control Markov dynamics with unobserved mechanism context:
\begin{align}
    s_t &\in \R^{d_s}, &
    a_t &\in \R^{d_a}, &
    z_t &\in \R^{d_z}, \\
    s_{t+1} &= f^\star(s_t,a_t,z_t) + \epsilon_t.
\end{align}
The context $z_t$ represents physical or transition-level perturbations such as friction, mass, damping, delay, external drift, sticky transitions, and impulses. During offline data collection in the synthetic MuJoCo benchmark, the context labels are known, but at planning time the model must act from history and current state only.

The target is not merely low average transition error. A useful world model should support planning:
\begin{equation}
    \pi_{\hat f}
    =
    \arg\max_{\pi}
    \E_{\hat f}
    \left[
        \sum_{t=0}^{H-1} r(s_t,a_t,s_{t+1})
    \right],
\end{equation}
and should transfer this planning behavior back to the real or perturbed environment. This motivates reporting both prediction metrics such as $R^2@1$ and $R^2@10$, and control metrics such as MPC return.

\section{Core Method}

\subsection{LLM as a constrained prior builder}

DUC-WM does not ask the LLM to write a simulator or arbitrary symbolic equations. Instead, the LLM must return JSON under a fixed schema:
\begin{equation}
    \mathcal{T}
    =
    \{T_j\}_{j=1}^K,
\end{equation}
where each mechanism template $T_j$ contains:
\begin{align}
    T_j =
    (&\mathrm{name}_j,\ I^s_j,\ I^a_j,\ I^o_j,\ \ell_j,\ g_j,\ \sigma_j,\ \rho_j,\ \tau_j).
\end{align}
Here $I^s_j$, $I^a_j$, and $I^o_j$ are sparse state, action, and output masks; $\ell_j$ is a safe DSL law type; $g_j$ is a small law gain; $\sigma_j$ is prior uncertainty; $\rho_j \in [0,1]$ is prior confidence; and $\tau_j$ marks whether the mechanism is slow, event-like, or unknown.

Allowed law types include:
\begin{equation}
\ell_j \in
\{\texttt{actuation},
\texttt{external\_drift},
\texttt{velocity\_damping},
\texttt{inertia\_shift},
\texttt{action\_delay},
\texttt{sticky\_velocity},
\texttt{impulse},
\texttt{gravity\_shift},
\texttt{learned\_residual}\}.
\end{equation}
The JSON is then parsed and compiled by code into a tensor-valued prior law:
\begin{equation}
    P_j(s_t,a_t,h_t) \in \R^{d_s}.
\end{equation}
The LLM prior is therefore inspectable, bounded, executable, and cannot inject arbitrary code.

\subsection{History-conditioned mechanism inference}

Let $h_t$ denote a short history window:
\begin{equation}
    h_t =
    (s_{t-L:t}, a_{t-L:t-1}).
\end{equation}
DUC-WM learns a context encoder:
\begin{align}
    q_\phi(u_t \mid h_t)
    &=
    \mathcal{N}(\mu_\phi(h_t), \diag(\sigma^2_\phi(h_t))), \\
    \alpha_t
    &=
    c \odot \tanh(u_t),
\end{align}
where $\alpha_{t,j}$ is the inferred strength of mechanism $j$, and $c$ is a per-mechanism scale derived from the template. At deployment, $\alpha_t$ is inferred from history; no oracle context is provided to DUC-WM.

\subsection{Mechanism-residual transition model}

The transition is modeled in delta form:
\begin{equation}
    \hat s_{t+1}
    =
    s_t
    +
    B_\theta(s_t,a_t)
    +
    C_\theta(s_t,a_t,\alpha_t)
    +
    \sum_{j=1}^{K}
    \alpha_{t,j}
    \left[
        \beta_j q_j P_j(s_t,a_t,h_t)
        +
        \gamma_t R_{\theta,j}(s_t,a_t)
    \right].
    \label{eq:duc-transition}
\end{equation}
The terms are:
\begin{itemize}[leftmargin=1.5em]
    \item $B_\theta$: shared base dynamics trunk.
    \item $C_\theta$: context-conditioned dynamics adapter.
    \item $P_j$: compiled prior law for mechanism $j$.
    \item $q_j \in [0,1]$: data-validated prior gate.
    \item $\beta_j > 0$: learned or calibrated prior strength.
    \item $R_{\theta,j}$: mechanism-specific neural residual.
    \item $\gamma_t \in [0,1]$: residual warmup gate during training.
\end{itemize}
This decomposition gives the model a route to use physics-like priors when they are helpful, while still retaining enough neural capacity to repair missing or incorrect priors.

The model also predicts diagonal transition uncertainty:
\begin{equation}
    p_\theta(s_{t+1} \mid s_t,a_t,h_t)
    =
    \mathcal{N}
    \left(
        \hat s_{t+1},
        \diag(\exp(\ell_\theta(s_t,a_t,\hat s_{t+1},\alpha_t)))
    \right).
\end{equation}

\section{Prior Validation}

LLM-generated priors may be partially correct, irrelevant, or wrong. DUC-WM validates each prior direction against observed transition deltas before training. For a candidate prior effect $P_j$ and target delta $\Delta s_t=s_{t+1}-s_t$, validation computes a weighted alignment over the template output dimensions $I^o_j$.

Let
\begin{align}
    p_{t,j} &= P_j(s_t,a_t,h_t)_{I^o_j}, \\
    y_t &= \Delta s_{t,I^o_j}.
\end{align}
The validation estimates a scalar alignment and least-squares scale:
\begin{align}
    \hat \lambda_j
    &=
    \frac{\sum_t \langle p_{t,j}, y_t\rangle_w}
         {\sum_t \|p_{t,j}\|^2_w + \varepsilon}, \\
    \mathrm{improve}_j
    &=
    \left[
    1 -
    \frac{\sum_t \|\hat\lambda_j p_{t,j}-y_t\|^2_w}
         {\sum_t \|y_t\|^2_w+\varepsilon}
    \right]_+, \\
    \mathrm{cos}_j
    &=
    \left[
    \frac{\sum_t \langle p_{t,j},y_t\rangle_w}
         {\sqrt{\sum_t\|p_{t,j}\|^2_w}\sqrt{\sum_t\|y_t\|^2_w}+\varepsilon}
    \right]_+.
\end{align}
The raw validation score is:
\begin{equation}
    v_j = \sqrt{\mathrm{improve}_j \cdot \mathrm{cos}_j}.
\end{equation}
It is converted into a gate:
\begin{equation}
    q_j =
    \begin{cases}
        0, & v_j \approx 0, \\
        q_{\min} + (1-q_{\min})\frac{v_j}{v_j+\tau}, & \text{otherwise}.
    \end{cases}
\end{equation}
The prior strength is initialized or updated as:
\begin{equation}
    \beta_j = \clip(|\hat\lambda_j|,\ \beta_{\min},\ \beta_{\max}).
\end{equation}

For the main supervised-context benchmark, prior validation may use context labels during training. For a stricter no-label variant, the implementation supports:
\begin{equation}
    \texttt{context\_weight}=0,\quad
    \texttt{teacher\_force\_context}=\texttt{false},\quad
    \texttt{prior\_validation\_use\_context}=\texttt{false},
\end{equation}
so validation uses only $(s_t,a_t,s_{t+1})$ and does not peek at true context labels.

\section{Training Objective}

The DUC-WM training objective combines predictive likelihood, latent regularization, optional context supervision, control-sensitive prediction loss, rollout consistency, and structure regularizers:
\begin{align}
    \mathcal{L}
    &=
    \mathcal{L}_{\mathrm{NLL}}
    +
    \lambda_{\mathrm{KL}}\mathcal{L}_{\mathrm{KL}}
    +
    \lambda_{\mathrm{ctx}}\mathcal{L}_{\mathrm{ctx}}
    +
    \lambda_{\mathrm{ctrl}}\mathcal{L}_{\mathrm{ctrl}}
    +
    \lambda_{\mathrm{roll}}\mathcal{L}_{\mathrm{roll}}
    \nonumber \\
    &\quad+
    \lambda_{\mathrm{res}}\mathcal{L}_{\mathrm{res}}
    +
    \lambda_{\mathrm{TR}}\mathcal{L}_{\mathrm{TR}}
    +
    \lambda_{\mathrm{orth}}\mathcal{L}_{\mathrm{orth}}
    +
    \lambda_{\mathrm{sparse}}\mathcal{L}_{\mathrm{sparse}}
    +
    \lambda_{\mathrm{unk}}\mathcal{L}_{\mathrm{unk}}
    +
    \lambda_{\beta}\mathcal{L}_{\beta}.
\end{align}

The likelihood term is:
\begin{equation}
    \mathcal{L}_{\mathrm{NLL}}
    =
    -\log p_\theta(s_{t+1}\mid s_t,a_t,h_t).
\end{equation}
The latent KL term regularizes the inferred mechanism strengths toward the template prior:
\begin{equation}
    \mathcal{L}_{\mathrm{KL}}
    =
    \KL
    \left(
        q_\phi(u_t\mid h_t)
        \,\Vert\,
        \mathcal{N}(\mu_0,\diag(\sigma_0^2))
    \right).
\end{equation}
When synthetic context labels are used, context supervision is:
\begin{equation}
    \mathcal{L}_{\mathrm{ctx}}
    =
    \|\alpha_t - z_t\|_2^2.
\end{equation}
This is matched by the main fair baseline CaDM-supervised; both DUC-WM and CaDM-supervised may learn from context labels during training but must infer context at planning time.

The control-weighted prediction loss emphasizes reward-sensitive dimensions:
\begin{equation}
    \mathcal{L}_{\mathrm{ctrl}}
    =
    \|W_r^{1/2}(\hat s_{t+1}-s_{t+1})\|_2^2.
\end{equation}
The multi-step rollout term is:
\begin{equation}
    \mathcal{L}_{\mathrm{roll}}
    =
    \sum_{\ell=1}^{H_r}
    \|W_r^{1/2}(\hat s_{t+\ell}-s_{t+\ell})\|_2^2.
\end{equation}

\subsection{Trust-region residual regularization}

For each mechanism, high-confidence priors should not be overwritten too freely, while low-confidence priors should allow more neural repair. Let:
\begin{align}
    \rho_j^{\mathrm{eff}}
    &=
    \rho_j \cdot d_j, \\
    \delta_j
    &=
    \delta_{\min} + (1-\rho_j^{\mathrm{eff}})\delta_{\mathrm{range}},
\end{align}
where $d_j$ is the data confidence produced by prior validation. The trust-region penalty is:
\begin{equation}
    \mathcal{L}_{\mathrm{TR}}
    =
    \frac{1}{K}
    \sum_{j=1}^{K}
    \rho_j^{\mathrm{eff}}
    \left[
        \|R_{\theta,j}\|_2
        -
        \delta_j\|P_j\|_2
    \right]_+^2.
\end{equation}
This encourages residuals to repair, rather than erase, validated priors.

\section{Planning}

All methods are evaluated with the same learned reward surrogate and CEM-MPC planner. A reward model $\hat r_\psi$ is trained from the same offline transitions:
\begin{equation}
    \hat r_\psi(s_t,a_t,s_{t+1})
    \approx
    r_t.
\end{equation}
CEM-MPC optimizes a horizon-$H$ action sequence under the learned dynamics:
\begin{equation}
    a_{t:t+H-1}^{\star}
    =
    \arg\max_{a_{t:t+H-1}}
    \sum_{\ell=0}^{H-1}
    \hat r_\psi
    \left(
        \hat s_{t+\ell},
        a_{t+\ell},
        \hat s_{t+\ell+1}
    \right).
\end{equation}
The first action is executed, and the plan is recomputed at the next step. The current main protocol sets the planner uncertainty penalty to zero:
\begin{equation}
    \lambda_{\mathrm{uncertainty}}=0,
\end{equation}
so the PETS-style ensemble baseline is not penalized for exposing variance while single-model baselines expose non-equivalent uncertainty estimates.

\section{Baselines and Fairness}

The main workshop comparison uses four methods:
\begin{center}
\begin{tabular}{lll}
\toprule
Method & Role & Fairness interpretation \\
\midrule
DUC-WM & Proposed method & Infers context from history at planning time \\
CaDM-supervised & Closest adaptive baseline & Uses context supervision during training, infers at planning \\
PETS-context & Oracle-context upper bound & Receives true context during evaluation/planning \\
CaDM-context & Oracle-context upper bound & Receives true context during evaluation/planning \\
\bottomrule
\end{tabular}
\end{center}

This comparison is intentionally conservative. CaDM-supervised is the closest fair competitor because it receives the same type of training-time context supervision as DUC-WM, but must infer context at planning time. PETS-context and CaDM-context are not fair in the sense of equal information; they are stronger oracle baselines. If DUC-WM matches or outperforms them in planning return, the result is especially meaningful.

The protocol also controls several common confounders:
\begin{itemize}[leftmargin=1.5em]
    \item shared train/test transition budgets for all methods;
    \item same random seed per environment/method comparison;
    \item same CEM-MPC planner;
    \item same learned reward surrogate protocol;
    \item reward model re-seeding to avoid RNG drift across models;
    \item no PETS uncertainty penalty in the main planner;
    \item DUC reward-wake, action-rank, and learned symbolic revalidation disabled in the main protocol.
\end{itemize}

\section{Experimental Protocol}

The current main benchmark uses five MuJoCo environments:
\begin{equation}
    \{\text{Swimmer},\text{Hopper},\text{Walker2d},\text{HalfCheetah},\text{InvertedDoublePendulum}\}.
\end{equation}
Pusher and Ant are implemented as appendix-scale stress tests.

The main run uses:
\begin{align}
    &\text{seeds} = \{0,1,2\}, \\
    &\text{methods} =
    \{\text{DUC-WM},\text{CaDM-supervised},\text{PETS-context},\text{CaDM-context}\}.
\end{align}
The core configuration is intentionally non-toy:
\begin{itemize}[leftmargin=1.5em]
    \item 40k--100k training transitions depending on the environment;
    \item 10k--20k test transitions;
    \item 80--100 training epochs;
    \item hidden size 384, three layers for most environments;
    \item PETS ensemble size 7;
    \item CEM-MPC with 512 candidates, 64 elites, and 4 iterations;
    \item 3 planning episodes per seed;
    \item BF16 training on GPU with CPU-parallel MuJoCo collection.
\end{itemize}

\section{Metrics and Final Table}

The final workshop table should report three quantities:
\begin{align}
    R^2@1,\quad R^2@10,\quad \text{MPC Reward}.
\end{align}
For horizon $H$, rollout $R^2$ is:
\begin{equation}
    R^2@H
    =
    1 -
    \frac{
        \sum_i
        \|s_{i+H}-\hat s_{i+H}\|_2^2
    }{
        \sum_i
        \|s_{i+H}-\bar s_H\|_2^2
        +
        \varepsilon
    }.
\end{equation}
The planning return is:
\begin{equation}
    J_{\mathrm{MPC}}
    =
    \E
    \left[
        \sum_{t=0}^{T-1}
        r(s_t,a_t,s_{t+1})
    \right],
\end{equation}
where actions are selected by CEM-MPC using the learned model.

The final ranking uses a composite average rank:
\begin{equation}
    \mathrm{AvgRank}(m)
    =
    \frac{1}{|\mathcal{E}||\mathcal{M}|}
    \sum_{e\in\mathcal{E}}
    \sum_{q\in\mathcal{M}}
    \mathrm{rank}_{e,q}(m),
\end{equation}
where:
\begin{equation}
    \mathcal{M}
    =
    \{R^2@1,\ R^2@10,\ J_{\mathrm{MPC}}\}.
\end{equation}
Lower AvgRank is better. This avoids overclaiming from a single metric and explicitly includes planning utility, which is the most relevant downstream metric for a model-based control framework.

\section{Novelty Assessment}

\subsection{What is not individually novel}

A strict reviewer will note that several ingredients are known:
\begin{itemize}[leftmargin=1.5em]
    \item probabilistic ensemble dynamics models are standard in model-based RL;
    \item latent-context dynamics models are known for adapting to hidden environment shifts;
    \item residual dynamics models are common;
    \item trust-region style regularization is a known training principle;
    \item LLM-generated priors are an emerging theme.
\end{itemize}
Therefore, the paper should not claim that any one of these ingredients alone is the main novelty.

\subsection{What is novel enough for a workshop}

The novelty of DUC-WM is the framework-level composition:
\begin{enumerate}[leftmargin=1.5em]
    \item \textbf{LLM-to-executable-prior interface.}
    The LLM produces constrained mechanism templates, not free-form code. The output is compiled into safe tensor laws.

    \item \textbf{Data-validated law priors.}
    The model validates each law against observed deltas and turns validation into gates and calibrated strengths, so wrong priors become weak hints rather than hard constraints.

    \item \textbf{Mechanism-specific residual repair.}
    Each named mechanism has a residual channel, so the neural model learns what the prior misses while preserving an interpretable decomposition.

    \item \textbf{Hidden-context inference for planning.}
    The model infers mechanism strengths from history and uses them inside MPC, rather than assuming oracle context at deployment.

    \item \textbf{Planning-aware evaluation.}
    The benchmark explicitly distinguishes prediction quality from model utility for MPC, using a composite AvgRank over $R^2@1$, $R^2@10$, and reward.
\end{enumerate}

This is sufficiently novel for a workshop if the paper frames the contribution as a \emph{framework and empirical study} rather than a completely new theoretical learning algorithm.

\section{Recommended Claims}

If the 5-environment benchmark produces the expected pattern, a defensible claim is:

\begin{quote}
DUC-WM achieves the best composite average rank across prediction accuracy and MPC return on hidden-mechanism MuJoCo benchmarks, while matching or improving planning return against context-supervised and oracle-context model-based baselines.
\end{quote}

The paper should avoid stronger claims such as:
\begin{itemize}[leftmargin=1.5em]
    \item ``DUC-WM universally outperforms PETS/CaDM.''
    \item ``The LLM prior alone is responsible for the improvement.''
    \item ``The learned mechanisms are fully causally identified.''
\end{itemize}

Current diagnostic metrics suggest that the residual path often carries much of the transition delta. Thus, a safer interpretation is:
\begin{quote}
The mechanism-prior structure improves the planning utility and organization of the world model, even when neural residuals remain necessary for accurate dynamics.
\end{quote}

\section{Limitations and Follow-Up Experiments}

\begin{enumerate}[leftmargin=1.5em]
    \item \textbf{Mechanism identifiability.}
    High reward does not guarantee that each inferred mechanism aligns perfectly with the true synthetic context.

    \item \textbf{LLM prior attribution.}
    To isolate the LLM prior, report ablations such as DUC-no-LLM, DUC-random-prior, and DUC-wrong-prior in the appendix.

    \item \textbf{No-label setting.}
    A stricter no-label DUC variant is implemented. It should be used as an appendix stress test, not the main workshop protocol.

    \item \textbf{Canonical PETS.}
    The current PETS baseline is a PETS-style probabilistic ensemble and does not implement full particle trajectory sampling inside MPC. The paper should call it PETS-style unless a full PETS-TS planner is added.

    \item \textbf{Statistical power.}
    Three seeds over five environments is reasonable for a workshop. Stronger claims would require more seeds and possibly more environments.
\end{enumerate}

\section{Bottom Line}

DUC-WM has enough novelty for a workshop submission if positioned correctly:
\begin{itemize}[leftmargin=1.5em]
    \item not as ``LLM writes physics'';
    \item not as a pure symbolic method;
    \item not as a universal replacement for PETS or CaDM;
    \item but as a controlled mechanism-residual world model that converts LLM semantic priors into validated executable priors, repairs them with neural residuals, and evaluates success through both prediction and planning.
\end{itemize}

The strongest expected empirical story is:
\begin{quote}
DUC-WM may not always dominate every $R^2$ column, but it should achieve the best or near-best planning return and the best composite AvgRank, showing that mechanism-structured world models can be more useful for control than black-box prediction accuracy alone suggests.
\end{quote}

\end{document}
